\newpage
% ==========================================
% ROZDZIAŁ 1
% ==========================================
\section{Teoretyczne podstawy interakcji człowiek-rój (HSI) oraz architektury agentowo-robotycznej}

\subsection{Bariery kognitywne i ograniczenia poznawcze w wielodronowych systemach sterowania}
Przejście od zdalnego sterowania (ang. \textit{teleoperation}) do sterowania nadzorczego (ang. \textit{supervisory control}) stanowi fundamentalny wymóg w rozwoju zaawansowanych systemów UAV. W tradycyjnym ujęciu, złożoność poznawcza sterowania niezależnymi robotami skaluje się liniowo ze wzorem $O(n)$, gdzie $n$ to liczba jednostek. Wymaga to od operatora alokacji zasobów uwagi wprost proporcjonalnie do wielkości floty, co w warunkach krytycznych prowadzi do drastycznego przeciążenia kognitywnego (ang. \textit{cognitive overload}). Zgodnie z badaniami przeprowadzonymi przez Kollinga i in., nadrzędnym celem inżynierii systemów rojowych jest sprowadzenie tej złożoności do poziomu $O(1)$, w którym operator zarządza całym rojem (lub jego podzbiorami -- ang. \textit{subswarms}) jako pojedynczą, zintegrowaną encją \cite{kolling2016}.

Pojemność ludzkiej pamięci roboczej w kontekście jednoczesnej kontroli wielokrotnych bezzałogowców jest ściśle ograniczona. Badania empiryczne wykazują, że ludzki operator jest w stanie skutecznie śledzić do 16 obiektów wyłącznie w przypadku najprostszych zadań selekcji, do 7 w operacjach o średniej złożoności, natomiast w przypadku heterogenicznych, zaawansowanych misji -- zaledwie do 4 jednostek \cite{dia2010}. Próba przekroczenia tych barier, określana problematyką wielozadaniowości (ang. \textit{multitasking}), skutkuje skokową degradacją świadomości sytuacyjnej (ang. \textit{Situation Awareness - SA}), w szczególności podczas tzw. przełączania zadań (ang. \textit{task switching}) \cite{chen2010}. W interakcji z zaawansowaną automatyzacją zjawisko to rodzi dodatkowe zagrożenia, takie jak \textit{cognitive tunneling} (zawężenie uwagi do jednego wskaźnika) oraz nadmierna ufność w system (ang. \textit{complacency}), która zwiększa ryzyko występowania błędów decyzyjnych (tzw. zjawiska \textit{false-alarm} i \textit{miss-prone}) \cite{chen2010}.

Rozwiązaniem tego problemu jest wdrożenie paradygmatu współpracy wielu operatorów z wieloma platformami (M:N) oraz architektury zespołu ludzko-autonomicznego (ang. \textit{Human-Autonomy Teaming, HAT}), wprowadzającej role takich jak monitory czy menedżerowie lotu \cite{smith2021}. Odpowiednio zaprojektowany interfejs w takim zespole pozwala na zastosowanie zarządzania przez wyjątki (ang. \textit{Management by Exception - MBE}) lub zgody (ang. \textit{Management by Consent - MBC}) \cite{chen2010}. Co więcej, wprowadzono koncepcję tzw. ,,dobroczynnego zaniedbania'' (ang. \textit{neglect benevolence}), która zakłada, że w systemach rojowych o wysokim wskaźniku autonomii, czasowa nieobecność interwencji operatora często pozwala rojom skuteczniej osiągać optymalne, globalne stany stabilności \cite{kolling2016}.

\subsection{Modele komunikacji głosowej i interfejsy HRI w standardzie ROS}
Aby zrealizować model sterowania na poziomie $O(1)$, konieczne jest zastąpienie klasycznych joysticków interfejsami o wysokim poziomie semantycznym. Interakcja Człowiek-Robot (ang. \textit{Human-Robot Interaction, HRI}) opiera się obecnie na dwukierunkowej komunikacji (ang. \textit{bidirectional communication}), która pozwala robotowi nie tylko na pasywny odbiór komend, lecz także na informowanie człowieka o swoich zamiarach (ang. \textit{intent expression}) \cite{pomianek2026}.

W nowoczesnych architekturach robotycznych dąży się do multimodalnej fuzji danych percepcyjnych. Integracja sygnałów fonicznych z estymacją postawy, mimiką czy wskaźnikami biologicznymi (np. sygnałami EEG) pozwala na zbudowanie wysoce odpornego modelu predykcyjnego stanów emocjonalnych i intencji operatora \cite{pomianek2026}. Ekosystem Robot Operating System (ROS) standaryzuje to podejście poprzez zbiór konwencji \texttt{ROS4HRI} (w szczególności \texttt{REP-155}). Umożliwia to zunifikowany przesył danych za pośrednictwem dedykowanych bibliotek takich jak \texttt{hri\_msgs} czy pakietów śledzących, m.in. \texttt{hri\_person\_manager} \cite{pomianek2026}. 

Krytycznym aspektem rozwoju systemów rozpoznawania mowy dla rojów dronów na polu walki i w warunkach infrastruktury krytycznej są kwestie etyki i bezpieczeństwa. W dobie rygorystycznych wymogów prawnych (m.in. RODO/GDPR), preferowanym wzorcem architektonicznym jest Przetwarzanie na Krawędzi (ang. \textit{Edge Computing}), które gwarantuje asynchroniczną ewaluację modeli ML (Machine Learning) bez konieczności przesyłania wrażliwych danych biometrycznych poza sieć lokalną roju. Podejście to, znane jako \textit{Privacy by Design}, minimalizuje ryzyko wektorów ataków cyfrowych na komunikację akustyczną \cite{pomianek2026}.

\subsection{Paradygmat Embodied AI i rola wielkich modeli językowych (LLM) w sterowaniu}
Tłumaczenie nieskodyfikowanej mowy ludzkiej na formalne wywołania węzłów ROS2 wymaga warstwy pośredniczącej, wyposażonej w szeroki aparat kognitywny. Włączenie Wielkich Modeli Językowych (LLM) otwiera nowe możliwości w zakresie podejmowania decyzji wieloagentowych (ang. \textit{Multi-Agent Decision-Making, MADM}). W modelu tym algorytmy LLM ewoluują od prostych interpreterów poleceń do funkcji \textit{policy generators} oraz \textit{critics}, stymulujących mechanizmy tzw. uczenia w kontekście (ang. \textit{in-context learning}) oraz implementacji obliczeniowej Teorii Umysłu (ang. \textit{Theory of Mind - ToM}) u agentów sztucznych \cite{sun2025}.

Zastosowanie LLM w robotyce wymusiło powstanie paradygmatu Sztucznej Inteligencji Ucieleśnionej (ang. \textit{Embodied AI}). Samo oprogramowanie LLM nie dysponuje fizyczną świadomością granic czasoprzestrzennych, stąd jego integracja w zamkniętej pętli wymaga dedykowanych architektur pośrednich. Jednym z wiodących rozwiązań w tej dziedzinie jest framework RAI \cite{rachwal2025}. Architektura RAI pozwala na kategoryzację systemów inteligentnych na agenty konwersacyjne (ang. \textit{Conversational Agents}) i oparte na stanach (ang. \textit{State-Based Agents}). Poprzez wbudowane mechanizmy integracyjne (ang. \textit{Connectors}), w tym zwłaszcza integrację warstwy \texttt{ROS2Connector}, framework ten pozwala na bezproblemową komunikację modeli generatywnych z komponentami sprzętowymi \cite{rachwal2025}.

Centralnym wyzwaniem pozostaje zjawisko tzw. halucynacji AI, czyli konfabulowania przez modele nielogicznych i z fizycznego punktu widzenia niewykonalnych instrukcji nawigacyjnych. Aby ograniczyć to zjawisko, ucieleśnienie odbywa się przy pomocy technik Generowania Wzbogaconego Wyszukiwaniem (ang. \textit{Retrieval-Augmented Generation, RAG}). Mechanizm ten wykorzystuje zaawansowane, wielowymiarowe bazy wektorowe (np. \texttt{Faiss}) do semi-automatycznego wstrzykiwania do \textit{system promptów} wiedzy kontekstowej o parametrach roju, limitach kinematycznych i dostępnych zachowaniach (ang. \textit{Tools}). Pozwala to LLM na deterministyczne przyporządkowywanie otwartych instrukcji głosowych do zamkniętego zbioru algorytmów operacyjnych drona \cite{rachwal2025}.

\subsection{Algorytmy inteligencji stadnej i optymalizacja zachowań rojowych}
Na najniższym poziomie abstrakcji układu decyzyjnego, polecenia przetworzone przez LLM zostają delegowane do systemu planowania trajektorii i alokacji celów przestrzennych roju. Biologicznie inspirowana Inteligencja Stadna (ang. \textit{Swarm Intelligence}) wykorzystuje samoorganizujące się jednostki z ograniczonymi receptorami do powstawania skomplikowanych emergentnych zachowań globalnych \cite{bonabeau1999, kennedy2001}. W naturze zachowania takie -- jak formowanie klucza w kształcie litery V przez ptaki migrujące (optymalizacja oporów aerodynamicznych) lub systemy podziału ról (ang. \textit{division of labor}) u insektów -- opierają się nierzadko na tzw. stygmergii, tj. niebezpośredniej komunikacji poprzez modyfikację fizycznego środowiska \cite{bonabeau1999, kennedy2001}.

W warunkach symulowanej autonomii roju, zadania re-alokacji i nawigacji dynamicznej modelowane są najczęściej w matematycznej domenie Dynamicznego Problemu Komiwojażera (ang. \textit{Dynamic Traveling Salesman Problem, DTSP}) \cite{strak2017}. Adaptacja do szybko zmieniających się warunków operacyjnych wymaga wysoce nieliniowych metaheurystyk. Do czołowych z nich należy algorytm Dyskretnej Optymalizacji Roju Cząstek (ang. \textit{Discrete Particle Swarm Optimization, DPSO}) \cite{strak2017}.

Mechanika algorytmu DPSO, stosowana w kontekście nawigacji rojowej UAV, wykorzystuje rozbudowane struktury sąsiedztwa i heterogeniczną adaptację parametrów, co chroni sub-roje przed przedwczesną zbieżnością w minimach lokalnych. Co więcej, wzbogacenie algorytmu DPSO o mechanizmy tzw. pamięci feromonowej potęguje efektywność przeszukiwania topologii grafu operacyjnego poprzez nakładanie kar i nagród dla eksplorowanych węzłów ścieżki \cite{strak2017}. W celu estymacji optymalnych ścieżek stosuje się zaawansowane relaksacje metryczne, w tym tzw. estymację odległości opartą o dolne ograniczenia \textit{Held-Karp} oraz $\alpha$-miarę \textit{Helsgauna}, co pozwala modelom LLM na deterministyczną ewaluację propozycji planu przed wysłaniem komend wykonawczych do węzłów środowiska ROS2 \cite{strak2017}.